\documentclass{article}
\usepackage{graphicx}

\title{First LaTeX Experience}
\author{İbrahimcan Kiracı}
\date{August 2026}

\begin{document}

\maketitle

\section{Introduction}
This is my very first LaTeX document.

In-line mathematical equations are written using single dollar signs, like $E = mc^2$.

For standalone, centered, and emphasized formulas, we use double dollar signs:
$$F = m \cdot a$$
$$E = mc^2$$

\section{Kinetic Energy}

The classical kinetic energy of an object is given by:
$$E_k = \frac{1}{2} \cdot m \cdot v^2$$

Where:
* $E_k$ is the kinetic energy,
* $m$ is the mass,
* $v$ is the velocity.

\section{Potential Energy}

Gravitational potential energy:
$$E_p = m \cdot g \cdot h$$

Einstein's mass-energy equivalence:
$$E = m \cdot c^2$$

\section{Mechanical Energy and Friction}

Frictional force (where $\mu$ is the friction coefficient and $N$ is the normal force):
$$f_k = \mu \cdot N$$

Thermal energy generated by friction over distance $d$:
$$E_{thermal} = f_k \cdot d$$

Total mechanical energy of the system:
$$E_{mech} = E_k + E_p$$

Conservation of energy in the presence of friction:
$$E_{total} = E_{mech} + E_{thermal}$$

\section{Essential Mathematical Operations}

1. Square Root:
$$\sqrt{x}$$

2. n-th Root:
$$\sqrt[4]{x}$$

3. Limit (as x approaches 0):
$$\lim_{x \to 0} f(x)$$

4. Derivative (Leibniz notation):
$$\frac{df}{dx}$$

5. Definite Integral:
$$\int_{a}^{b} x^2 \, dx$$

6. Summation (Sigma notation):
$$\sum_{i=1}^{n} i$$

7. Trigonometric and Inverse Functions

$$\sin(\theta)^2 + \cos(\theta)^2 = 1$$

$$\tan(\theta) = \frac{\sin(\theta)}{\cos(\theta)}$$

$$\cot(\theta) = \frac{1}{\tan(\theta)}$$

$$\sec(x) = \frac{1}{\cos(x)}, \quad \csc(x) = \frac{1}{\sin(x)}$$

8. Inverse Trigonometric Functions (Arc):
$$\arctan(x) = \theta$$

$$\arcsin(x)^2 + \arccos(x)^2 = \frac{\pi}{2}$$

$$\sin(\theta)^2 + \cos(\theta)^2 = 1$$

9. Logarithmic Functions

Natural logarithm:
$$\ln(x)$$

Logarithm with base 10:
$$\log_{10}(x)$$

\end{document}
